\problem Let $x$ be a nilpotent element of a ring $A$. Show that $1 + x$ is a unit of $A$.
Deduce that the sum of a nilpotent element and a unit is a unit.

\begin{proof}
    Let $x \in A$ be nilpotent.
    Then, $x \in \nilradical \subset \jacobsonradical$, so for all $y \in A$, $1- xy$ is a unit in $A$.
    In particular, $1 + x = 1 - x(-1)$ is a unit in $A$, as claimed.

    Let now $y$ be a unit of $A$.
    Then, $x + y = y(y^{-1}x + 1)$.
    As $x \in \nilradical$, $y^{-1}x \in \nilradical$ so that by the above, $y^{-1}x + 1$ is a unit in $A$.
    It follows that as a product of units in $A$, $x + y$ is also a unit in $A$.
\end{proof}

\setcounter{problem}{5}

\problem A ring $A$ is such that every ideal not contained in the nilradical contains a nonzero idempotent (that is, an element $e$ such that $e^2 = e \not= 0$). Prove that the nilradical and Jacobson radical of $A$ are equal.
\begin{proof}
    The Jacobson radical $\jacobsonradical$ of $A$ is an ideal of $A$ containing $\nilradical$.
    If the containment were proper, then by assumption, there is a nonzero element $e \in \jacobsonradical$ such that $e^2 = e$.
    Then, $e(e-1) = 0$. But, $e$ being in $\jacobsonradical$ means that $e-1$ is a unit in $A$, whence $e(e-1) = 0$ implies $e = 0$.
    Contradiction!
\end{proof}

\problem Let $A$ be a ring in which every element $x$ satisfies $x^n = x$ for some $n > 1$ (depending on $x$).
Show that every prime ideal in $A$ is maximal.

\begin{proof}
    Let $\frak{p} \ideal A$ be prime.
    We will show that the ring $A/\frak{p}$ is a field.
    Indeed, as $\frak{p}$ is prime, $A/\frak{p}$ is already an integral domain.
    Thus, it suffices to show that every nonzero $x + \frak{p}$ in $A/\frak{p}$ is invertible.
    To this end, let $x + \frak{p} \in A/\frak{p}$, not zero, be arbitrary.
    By assumption, there is an integer $n > 1$ so that $x^n = x$, or $x(x^{n-1} - 1) = 0$.
    Hence, $(x + \frak{p})(x^{n-1} -1  + \frak{p}) = \frak{p}$.
    Since $x + \frak{p}$ is not zero and $A/\frak{p}$ is an ID, it follows that $x^{n-1} - 1 + \frak{p} = \frak{p}$.
    Hence, $x^{n-1} + \frak{p} = 1 + \frak{p}$ so that $x + \frak{p}$ is a invertible.
    It follows that $A/\frak{p}$ is a field, whence $\frak{p} \ideal A$ is maximal.
\end{proof}

\problem Let $A$ be a ring $\not=\zero$. Show that the set of prime ideals of $A$ has minimal elements with respect to inclusion.
\begin{proof}
    Let $\Sigma$ be the set of prime ideals of $A$ ordered by inclusion.
    We will use Zorn's Lemma to show that every chain in $\Sigma$ has a lower bound, whence $\Sigma$ has a minimal element.
    Indeed, $\Sigma$ is nonempty since (by Zorn's Lemma) $A$ contains a maximal ideal.
    Let now $\{\frak{p}_\iota\}_{\iota \in I}$ be a chain in $\Sigma$ indexed by the set $I$ and let $\frak{p} = \bigcap_{\iota \in I} \frak{p}_\iota$.
    We claim that $\frak{p} \in \Sigma$.
    First, that $\frak{p}$ is an ideal of $A$ follows from $\frak{p}$ being an intersection of ideals.
    To show that $\frak{p}$ is prime, let $xy \in \frak{p}$ be arbitrary.
    Then, $xy \in \frak{p}_\iota$ for all $\iota \in I$.
    Suppose that $x\notin \frak{p}$ and $y \notin \frak{p}$.
    Then there are indices $\alpha, \beta \in I$ such that $x \notin \frak{p}_\alpha$ and $y\notin \frak{p}_\beta$.
    Since $\{\frak{p}_\iota\}$ is a chain, either $\frak{p}_\alpha \subseteq \frak{p}_\beta$ or $\frak{p}_\beta \subseteq \frak{p}_\alpha$.
    In either case, the smaller ideal contains neither $x$ nor $y$, but does contain $xy$. Contradiction!
    Thus, either $x\in \frak{p}$ or $y \in \frak{p}$. Hence, $\frak{p} \in \Sigma$.

    Of course, $\frak{p} \subseteq \frak{p}_\iota$ for every $\iota \in I$ so that $\frak{p}$ is a lower bound for $\{\frak{p}_\iota\}$.
    It follows by Zorn's Lemma that $\Sigma$ contains a minimal element, as claimed.
\end{proof}